% Part: many-valued-logic
% Chapter: three-valued-logics
% Section: lukasiewicz

\documentclass[../../../include/open-logic-section]{subfiles}

\begin{document}

\olfileid[fr]{mvl}{thr}{luk}

\olsection{La logique de \L ukasiewicz}

L'une des premières propositions de
logique plurivalente élaborées et
publiées est due au philosophe
polonais Jan \L ukasiewicz, en 1920.\footnote{\emph{O logice trójwartościowej},
\href{https://kpbc.umk.pl/dlibra/publication/250664/edition/247342}{\emph{Ruch Filozoficzny}, tome~5, no~9}, p.~170--171.}
Il s'inspirait du problème de la
bataille navale d'Aristote : il semble
que, \emph{aujourd'hui}, la phrase
«~Il y aura une bataille navale demain~»
ne soit ni vraie ni fausse; sa valeur
de vérité n'est pas encore fixée.
\L ukasiewicz proposa d'attribuer
une troisième valeur de vérité à ces
énoncés portant sur les «~futurs
contingents~».
\begin{quote}
Je peux supposer sans contradiction
que ma présence à Varsovie à un certain
moment de l'année prochaine, par exemple
à midi le 21 décembre, n'est à présent
déterminée ni positivement ni négativement.
Il est donc possible, mais non nécessaire,
que je sois à Varsovie à ce moment-là.
Selon cette hypothèse, la proposition
«~Je serai à Varsovie à midi le 21 décembre
de l'année prochaine~» ne peut, à présent,
être ni vraie ni fausse. Car, si elle
était vraie maintenant, ma présence
future à Varsovie devrait être nécessaire,
ce qui contredit l'hypothèse. Si elle
était fausse maintenant, en revanche,
ma présence future à Varsovie devrait
être impossible, ce qui contredit
également l'hypothèse. La proposition
considérée n'est donc pour le moment
ni vraie ni fausse et doit posséder
une troisième valeur, différente de
«~0~», ou le faux, et de «~1~», ou le
vrai. Nous pouvons désigner cette
valeur par «~$\frac{1}{2}$~».
Elle représente «~le possible~» et
s'ajoute au vrai et au faux comme
troisième valeur.
\end{quote}
Nous emploierons $\Undef$ pour la
troisième valeur de vérité de
\L ukasiewicz.\footnote{\L ukasiewicz
emploie ici «~possible~» dans un sens
inhabituel aujourd'hui : possible,
mais non nécessaire.}

Les fonctions de vérité des connecteurs
$\lnot$, $\land$ et $\lor$ se
déterminent facilement selon cette
interprétation. La négation d'un
énoncé portant sur un futur contingent
porte elle aussi sur un futur contingent;
ainsi, $\tf{\lnot}(\Undef)=\Undef$.
Si l'un des membres d'une conjonction
est indéterminé et l'autre vrai,
la conjonction est elle aussi
indéterminée : selon l'issue du futur
contingent, elle pourrait devenir
vraie ou fausse. Nous posons donc
\[
  \tf{\land}(\True,\Undef)
  =\tf{\land}(\Undef,\True)
  =\Undef.
\]
En revanche, si l'autre membre est
faux, la conjonction ne pourra pas
devenir vraie; ainsi,
\[
  \tf{\land}(\False,\Undef)
  =\tf{\land}(\Undef,\False)
  =\False.
\]
Lorsque les arguments ont déjà une
valeur fixée, $\True$ ou $\False$,
les autres valeurs coïncident avec
celles de la logique classique.

Pour le conditionnel, la situation
est un peu plus délicate. Supposons
que $q$ porte sur un futur contingent.
Si $p$ est faux, alors $p\lif q$
sera vrai quelle que soit l'issue
de $q$; nous posons donc
$\tf{\lif}(\False,\Undef)=\True$.
Si $p$ est vrai, alors $q\lif p$
sera vrai quelle que soit l'issue
de $q$, d'où
$\tf{\lif}(\Undef,\True)=\True$.
Si $p$ est vrai, $p\lif q$ pourra
devenir vrai ou faux, d'où
$\tf{\lif}(\True,\Undef)=\Undef$.
De même, si $p$ est faux, $q\lif p$
pourra devenir vrai ou faux, d'où
$\tf{\lif}(\Undef,\False)=\Undef$.
Reste le cas où $p$ et $q$ portent
tous deux sur des futurs contingents.
La motivation initiale nous conduirait
à attribuer $\Undef$ à ce cas. Mais
alors $!A\lif !A$ ne serait \emph{pas}
une tautologie. \L ukasiewicz acceptait
de renoncer à $!A\lor\lnot !A$ et à
$\lnot(!A\land\lnot !A)$, mais non
à $!A\lif !A$. Il posa donc
$\tf{\lif}(\Undef,\Undef)=\True$.

\begin{defn}\ollabel{def:lukasiewicz}
La logique à trois valeurs de
\L ukasiewicz est définie par la matrice
suivante :
\begin{enumerate}
  \item Le fragment de $\Lang L_0$ dont les connecteurs sont $\lnot$,
  $\land$, $\lor$ et $\lif$.
  \item L'ensemble des valeurs de vérité
  $V=\{\True,\Undef,\False\}$.
  \item $\True$ est la seule valeur
  désignée, c'est-à-dire
  $V^+=\{\True\}$.
  \item Les fonctions de vérité sont
  données par les tables suivantes :
  \begin{center}
    \begin{tabular}{c|c}
      $\tf{\lnot}$ & \\
      \hline
      $\True$ & $\False$ \\
      $\Undef$ & $\Undef$ \\
      $\False$ & $\True$
    \end{tabular}
    \quad
    \begin{tabular}{c|ccc}
      $\tf{\land}[\LogLuk[3]]$ & $\True$ & $\Undef$ & $\False$ \\
      \hline
      $\True$ & $\True$ & $\Undef$ & $\False$ \\
      $\Undef$ & $\Undef$ & $\Undef$ & $\False$ \\
      $\False$ & $\False$ & $\False$ & $\False$
    \end{tabular}
    \\[2ex]
    \begin{tabular}{c|ccc}
      $\tf{\lor}[\LogLuk[3]]$ & $\True$ & $\Undef$ & $\False$ \\
      \hline
      $\True$ & $\True$ & $\True$ & $\True$ \\
      $\Undef$ & $\True$ & $\Undef$ & $\Undef$ \\
      $\False$ & $\True$ & $\Undef$ & $\False$
    \end{tabular}
    \quad
    \begin{tabular}{c|ccc}
      $\tf{\lif}[\LogLuk[3]]$ & $\True$ & $\Undef$ & $\False$ \\
      \hline
      $\True$ & $\True$ & $\Undef$ & $\False$ \\
      $\Undef$ & $\True$ & $\True$ & $\Undef$ \\
      $\False$ & $\True$ & $\True$ & $\True$
    \end{tabular}
  \end{center}
\end{enumerate}
\end{defn}

Toute formule~$!A$ ne contenant
que $\lnot$, $\land$ et $\lor$
prend la valeur~$\Undef$ si toutes
ses variables propositionnelles
reçoivent~$\Undef$. Ainsi, les
tautologies classiques
$p\lor\lnot p$ et
$\lnot(p\land\lnot p)$ ne sont
pas des tautologies de~$\LogLuk[3]$ :
pour chacune d'elles,
$\pValue{v}(!A)=\Undef$ lorsque
$\pAssign{v}(p)=\Undef$.

Lorsque toutes les variables
propositionnelles de~$!A$ reçoivent
$\True$ ou $\False$, sa valeur
coïncide avec la valeur classique.

\begin{prop}
  Si $\pAssign{v}(p)\in
  \{\True,\False\}$ pour toute
  variable propositionnelle~$p$
  de~$!A$, alors
  $\pValue{v}(!A)[\LogLuk[3]]
  =\pValue{v}(!A)[\LogCL]$.
\end{prop}

\begin{prob}\label{mvl:thr:luk:prob:luk-iff}
  Supposons que nous définissions
  $\pValue{v}(!A\liff !B)
  =\pValue{v}((!A\lif !B)\land
  (!B\lif !A))$ dans~$\LogLuk[3]$.
  Quelle serait la table de vérité
  de~$\liff$?
\end{prob}

De nombreuses tautologies classiques
\emph{sont} aussi des tautologies de
$\LogLuk[3]$, par exemple
$\lnot p\lif(p\lif q)$. Comme en
logique classique, nous pouvons le
vérifier par une table de vérité :
\begin{center}
\begin{tabular}{cc|cccccc}
  $p$ & $q$ & $\lnot$ & $p$ & $\lif$ & $\smash{(}p$ & $\lif$ & $q\smash{)}$ \\ \hline
  \True & \True & \False & \True & \True & \True & \True & \True \\
  \True & \Undef & \False & \True & \True & \True & \Undef & \Undef \\
  \True & \False & \False & \True & \True & \True & \False & \False \\
  \Undef & \True & \Undef & \Undef & \True & \Undef & \True & \True \\
  \Undef & \Undef & \Undef & \Undef & \True & \Undef & \True & \Undef \\
  \Undef & \False & \Undef & \Undef & \True & \Undef & \Undef & \False \\
  \False & \True & \True & \False & \True & \False & \True & \True \\
  \False & \Undef & \True & \False & \True & \False & \True & \Undef \\
  \False & \False & \True & \False & \True & \False & \True & \False \\
\end{tabular}
\end{center}

\begin{prob}
  Montrer que les formules suivantes
  sont des tautologies de~$\LogLuk[3]$ :
  \begin{enumerate}
    \item $p\lif(q\lif p)$;
    \item\label{mvl:thr:luk:prob:luk-taut-2}
      $\lnot(p\land q)\liff
      (\lnot p\lor\lnot q)$;
    \item\label{mvl:thr:luk:prob:luk-taut-3}
      $\lnot(p\lor q)\liff
      (\lnot p\land\lnot q)$.
  \end{enumerate}
  Dans \olref[mvl][thr][luk]{prob:luk-taut-2}
  et \olref[mvl][thr][luk]{prob:luk-taut-3},
  prendre $!A\liff !B$ comme abréviation
  de $(!A\lif !B)\land(!B\lif !A)$,
  ou utiliser la réponse à
  l'\olref[mvl][thr][luk]{prob:luk-iff}.
\end{prob}

\begin{prob}
  Montrer que les tautologies classiques
  suivantes ne sont pas des tautologies
  de~$\LogLuk[3]$ :
  \begin{enumerate}
    \item $(\lnot p\land p)\lif q$;
    \item $((p\lif q)\lif p)\lif p$;
    \item $(p\lif(p\lif q))\lif(p\lif q)$.
  \end{enumerate}
\end{prob}

On pourrait alors penser que, même si
toutes les tautologies classiques ne
sont pas des tautologies de~$\LogLuk[3]$,
elles prennent au moins la valeur
$\True$ ou $\Undef$ dans toute
valuation. Ce n'est pas le cas.
La formule
\[
  \lnot(p\lif\lnot p)
  \lor\lnot(\lnot p\lif p)
\]
en fournit un contre-exemple :
elle vaut $\False$ lorsque
$p$ vaut~$\Undef$.

\begin{prob}
  Parmi les relations suivantes,
  lesquelles valent dans la logique
  de~\L ukasiewicz? Donner une table
  de vérité pour chacune.
  \begin{enumerate}
    \item $p,p\lif q\Entails q$;
    \item $\lnot\lnot p\Entails p$;
    \item $p\land q\Entails p$;
    \item $p\Entails p\land p$;
    \item $p\Entails p\lor q$.
  \end{enumerate}
\end{prob}

\L ukasiewicz espérait construire
une logique de la possibilité à partir
de son système à trois valeurs, en
introduisant un connecteur unaire
$\Diamond !A$ (pour «~$!A$ est possible~»)
et un connecteur correspondant
$\Box !A$ (pour «~$!A$ est nécessaire~») :
\begin{center}
  \begin{tabular}{c|c}
    $\tf{\Diamond}$ & \\
    \hline
    $\True$ & $\True$ \\
    $\Undef$ & $\True$ \\
    $\False$ & $\False$
  \end{tabular}
  \quad
  \begin{tabular}{c|c}
    $\tf{\Box}$ & \\
    \hline
    $\True$ & $\True$ \\
    $\Undef$ & $\False$ \\
    $\False$ & $\False$
  \end{tabular}
\end{center}
Autrement dit, $p$ est possible si
et seulement si sa fausseté n'est
pas déjà établie; $p$ est nécessaire
si et seulement si sa vérité est déjà
établie.

\begin{prob}
  Montrer que
  $\Box p\liff\lnot\Diamond\lnot p$
  et $\Diamond p\liff\lnot\Box\lnot p$
  sont des tautologies de~$\LogLuk[3]$
  étendue par les tables de vérité de
  $\Box$ et de~$\Diamond$.
\end{prob}

Les limites de cette logique modale
proposée apparurent cependant rapidement.
Quelle que soit l'issue des événements,
$p\land\lnot p$ ne pourra jamais
devenir vrai. Même si sa valeur n'est
pas encore fixée et qu'il est donc
indéterminé, cet énoncé devrait être
impossible : $\lnot\Diamond
(p\land\lnot p)$ devrait donc être
une tautologie. Pourtant, si
$\pAssign{v}(p)=\Undef$, alors
$\pValue{v}(\lnot\Diamond
(p\land\lnot p))=\False$. La conjonction vaut alors $\Undef$,
sa possibilité vaut $\True$ et la négation de celle-ci vaut $\False$.
Le codage vérifonctionnel proposé ne rend donc pas compte de
l'impossibilité d'une contradiction, même avec cette troisième
valeur de vérité.

\end{document}
